\documentclass[11pt,a4paper]{article}
\usepackage[a4paper,margin=22mm]{geometry}
\usepackage{fontspec}
\usepackage{amsmath,mathtools}
\usepackage{unicode-math}
\usepackage{array,booktabs}
\usepackage{enumitem}
\setmainfont{DejaVu Serif}
\setmathfont{DejaVu Math TeX Gyre}
\pagestyle{plain}
\setlength{\parindent}{1.25em}
\setlength{\parskip}{0.35em}

\begin{document}

\begin{center}
{\Large\bfseries Нативный PDF: формулы математического анализа}\\[4pt]
{\small Формулы набраны LaTeX и не растеризованы}
\end{center}

Пусть $f$ непрерывна на отрезке $[a,b]$, а $F'(x)=f(x)$. Тогда формула Ньютона--Лейбница имеет вид
\begin{equation}
  \int_a^b f(x)\,dx = F(b)-F(a).
  \tag{N.1}
\end{equation}
Для функции двух переменных полный дифференциал и квадратичная форма второго порядка записываются так:
\begin{align}
  du &= \frac{\partial u}{\partial x}\,dx + \frac{\partial u}{\partial y}\,dy, \tag{N.2}\\
  d^2u &= u_{xx}\,dx^2 + 2u_{xy}\,dx\,dy + u_{yy}\,dy^2. \tag{N.3}
\end{align}
Ряд Фурье функции периода $2\pi$ задаётся формулой
\begin{equation}
 f(x) \sim \frac{a_0}{2}+\sum_{n=1}^{\infty}\bigl(a_n\cos nx+b_n\sin nx\bigr),
 \quad a_n=\frac{1}{\pi}\int_{-\pi}^{\pi}f(x)\cos nx\,dx.
 \tag{N.4}
\end{equation}
Для проверки распознавания пределов включим составную дробь:
\begin{equation}
 \lim_{x\to0}\frac{\sqrt{1+x}-1}{x}=\frac12,
 \qquad
 \lim_{n\to\infty}\left(1+\frac{x}{n}\right)^n=e^x.
 \tag{N.5}
\end{equation}

\newpage
\begin{center}
{\Large\bfseries Нативный PDF: линейная алгебра и системы}
\end{center}

Система линейных уравнений $A\mathbf{x}=\mathbf{b}$ имеет единственное решение, если $\det A\ne0$. В развёрнутом виде:
\begin{equation}
\begin{cases}
  2x-y+3z=7,\\
  x+4y-z=2,\\
  3x+2y+2z=9.
\end{cases}
\tag{N.6}
\end{equation}
Определитель и обратная матрица:
\begin{equation}
 \det A=\sum_{\sigma\in S_n}\operatorname{sgn}(\sigma)
 \prod_{i=1}^{n}a_{i,\sigma(i)},
 \qquad A^{-1}=\frac{1}{\det A}\operatorname{adj}A.
 \tag{N.7}
\end{equation}
Спектральное разложение симметрической матрицы:
\begin{equation}
 A=Q\Lambda Q^{\mathsf T},\qquad Q^{\mathsf T}Q=I,
 \qquad \Lambda=\operatorname{diag}(\lambda_1,\ldots,\lambda_n).
 \tag{N.8}
\end{equation}
Норма и скалярное произведение:
\begin{equation}
 \|\mathbf{x}\|_2=\sqrt{\sum_{i=1}^{n}|x_i|^2},
 \qquad \langle\mathbf{x},\mathbf{y}\rangle=\mathbf{x}^{\mathsf T}\mathbf{y}.
 \tag{N.9}
\end{equation}

\newpage
\begin{center}
{\Large\bfseries Нативный PDF: вероятность и статистика}
\end{center}

Для дискретной случайной величины $X$ математическое ожидание и дисперсия равны
\begin{align}
 \mathbb{E}X &= \sum_k x_k p_k, \tag{N.10}\\
 \operatorname{Var}(X) &= \mathbb{E}\bigl[(X-\mathbb{E}X)^2\bigr]
 =\mathbb{E}X^2-(\mathbb{E}X)^2. \tag{N.11}
\end{align}
Плотность нормального распределения:
\begin{equation}
 p(x)=\frac{1}{\sigma\sqrt{2\pi}}
 \exp\!\left[-\frac{(x-\mu)^2}{2\sigma^2}\right].
 \tag{N.12}
\end{equation}
Формула Байеса для полной группы гипотез $H_1,\ldots,H_m$:
\begin{equation}
 \mathbb{P}(H_i\mid A)=
 \frac{\mathbb{P}(H_i)\,\mathbb{P}(A\mid H_i)}
 {\sum_{j=1}^{m}\mathbb{P}(H_j)\,\mathbb{P}(A\mid H_j)}.
 \tag{N.13}
\end{equation}
В качестве теста комбинаторной записи используем биномиальное тождество
\begin{equation}
 \sum_{k=0}^{n}\binom{n}{k}p^k(1-p)^{n-k}=1.
 \tag{N.14}
\end{equation}

\end{document}
